% Shared student/instructor body; compile a wrapper, not this file directly.
\providecommand{\dd}{\,\mathrm{d}}
\AtBeginDocument{%
\setlength{\abovedisplayskip}{4pt plus 1pt minus 1pt}
\setlength{\belowdisplayskip}{4pt plus 1pt minus 1pt}
\setlength{\abovedisplayshortskip}{3pt plus 1pt}
\setlength{\belowdisplayshortskip}{3pt plus 1pt}
\setlength{\jot}{2pt}
}
\tcbset{handout base/.append style={before skip=2.5pt,after skip=2.5pt,before upper={\RaggedRight}}}

\HandoutSetup{NE 630}{12}{Effective Cross Sections and $k_{\infty}$}
  {FNRP Ch.~3 (effective cross sections); review Lessons 10--11}
\begin{document}
\MakeHandoutTitle
\textcolor{KSUDeepPurple}{\textbf{Primary objective.}}
Compute spectrum-averaged cross sections and reaction rates to estimate
$k_{\infty}$ for an infinite, homogeneous system.
\par\vspace{5pt}

\begin{handoutcolumns}
\HandoutSection{1}{Balance over all energies}
For steady state, no leakage, and scattering that conserves neutron number,
\begin{align*}
 \Sigma_t(E)\phi(E)
 &=\int_0^\infty\Sigma_s(E'\!\to E)\phi(E')\dd E'\\
 &\quad+\chi(E)s_f'''+s_{\mathrm{ext}}'''(E),\\
 s_f'''&=\int_0^\infty\nu(E)\Sigma_f(E)\phi(E)\dd E.
\end{align*}
Here $\int_0^\infty\chi(E)\dd E=1$ and
$\int_0^\infty\Sigma_s(E'\!\to E)\dd E=\Sigma_s(E')$.
Define reaction-rate densities and the total external source:
\[
 R_x=\int_0^\infty\Sigma_x(E)\phi(E)\dd E,
 \qquad s=\int_0^\infty s_{\mathrm{ext}}'''(E)\dd E.
\]
\begin{boardbox}{Integrate, then cancel scattering}
\[
 R_t=R_s+s_f'''+s
 \quad\Longrightarrow\quad
 \boxed{R_a=\MathBlank[0.88in]{s_f'''+s}}.
\]
All terms are rates per unit volume, in $\mathrm{cm^{-3}\,s^{-1}}$.
\end{boardbox}
\begin{checkpoint}
Why does scattering disappear from the all-energy balance,
but not from the balance at a particular energy?
\RuledLines{2}
\end{checkpoint}

\HandoutSection{2}{Preserve the reaction rate}
Let $\phi_{\mathrm{tot}}=\int_0^\infty\phi(E)\dd E$.
The \textbf{effective cross section} is defined by
\begin{fixedbox}{Flux-weighted average (FNRP 3.38)}
\[
 \boxed{\bar\Sigma_x=
 \frac{\int_0^\infty\Sigma_x(E)\phi(E)\dd E}
      {\phi_{\mathrm{tot}}}},
 \qquad R_x=\bar\Sigma_x\phi_{\mathrm{tot}}.
\]
\end{fixedbox}
Thus $\bar\Sigma_x$ is the expected value of $\Sigma_x(E)$
under $w(E)=\phi(E)/\phi_{\mathrm{tot}}$.
Multiplying the spectrum by a constant does not change the average.
\[
 \overline{\nu\Sigma_f}=
 \frac{\int_0^\infty\nu(E)\Sigma_f(E)\phi(E)\dd E}
      {\phi_{\mathrm{tot}}}.
\]
Average the \emph{product}; in general this is not
$\bar\nu\,\bar\Sigma_f$ for two independently flux-weighted averages.
\begin{boardbox}{Nonmultiplying, source-driven limit}
\[
 \bar\Sigma_a\phi_{\mathrm{tot}}=s
 \quad\Longrightarrow\quad
 \phi_{\mathrm{tot}}=\MathBlank[0.85in]{s/\bar\Sigma_a}.
\]
Check: $\bar\Sigma_a$ has units $\mathrm{cm^{-1}}$ and
$\phi_{\mathrm{tot}}$ has units $\mathrm{cm^{-2}\,s^{-1}}$.
\end{boardbox}

\switchcolumn
\RaggedRight
\HandoutSection{3}{Normalize the spectral shapes}
Write $\theta=k_{\mathrm B}T$ and let $j=T,I,F$ denote the
thermal, intermediate, and fast ranges. All bounds below are in eV.
\begin{center}
\renewcommand{\arraystretch}{1.5}
\begin{tabular}{@{}cll@{}}
\toprule
$g$ &  Interval $[E^{L}_g, E^{U}_g]$ & Shape $h_g(E)$\\
\midrule
$T$ & $[10^{-3},1]$ & $E e^{-E/\theta}$\\
$I$ & $[1,10^5]$ & $1/[\Sigma_t(E)E]$\\
$F$ & $[10^5,10^7]$ & $\chi(E)/\Sigma_t(E)$\\
\bottomrule
\end{tabular}
\end{center}
These are the Maxwellian, narrow-resonance, and approximate
fast-spectrum shapes from earlier lessons, \emph{not} exact spectra.
For this finite-range approximation, neglect the tails outside these bounds.  Then
\begin{align*}
 C_g^{-1}&=\int_{E^{L}_g}^{E^{U}_g}h_g(E)\dd E, &
 w_g(E)&=C_gh_g(E),\\
 \int_{E^{L}_g}^{E^{U}_g}w_g(E)\dd E&=1, &
 \phi_g&=\int_{E^{L}_g}^{E^{U}_g}\phi(E)\dd E.
\end{align*}
Then, within range $g$,
\[
 \boxed{\phi(E)\approx\phi_g w_g(E)}.
\]
\begin{checkpoint}[Shape versus amount]
Does finding $C_T$, $C_I$, and $C_F$ determine
$\phi_T$, $\phi_I$, and $\phi_F$? What information is still missing?
\RuledLines{3}
\end{checkpoint}

\HandoutSection{4}{Combine the range averages}
Using a normalized shape in each interval gives
\[
 \bar\Sigma_{x,j}\approx
 \int_{E^{L}_g}^{E^{U}_g}\Sigma_x(E)w_g(E)\dd E.
\]
\begin{boardbox}{Add rates before dividing by total flux}
\begin{align*}
 R_x&\approx\bar\Sigma_{x,T}\phi_T
       +\bar\Sigma_{x,I}\phi_I+\bar\Sigma_{x,F}\phi_F,\\
 \phi_{\mathrm{tot}}&\approx\phi_T+\phi_I+\phi_F.
\end{align*}
\[
 \boxed{\bar\Sigma_x\approx
 \frac{\bar\Sigma_{x,T}\phi_T+\bar\Sigma_{x,I}\phi_I
       +\bar\Sigma_{x,F}\phi_F}{\phi_T+\phi_I+\phi_F}}.
\]
\end{boardbox}
The weights are \textbf{fractions of the total flux}, not fractions
of the energy interval. If $\phi_T=\phi_I=\phi_F$, then
\[
 \bar\Sigma_x\approx\MathBlank[1.95in]{
 (\bar\Sigma_{x,T}+\bar\Sigma_{x,I}+\bar\Sigma_{x,F})/3}.
\]
\end{handoutcolumns}

\newpage
\begin{handoutcolumns}
\HandoutSection{5}{Three useful spectrum averages}
\textbf{Fission-spectrum average.}
If $\Sigma_t$ is approximately constant over the fast range,
\[
 \bar\sigma_{x,F}\approx
 \frac{\int_{\mathcal E_F}\sigma_x(E)\chi(E)\dd E}
      {\int_{\mathcal E_F}\chi(E)\dd E}.
\]
The denominator is one only when the chosen interval contains
all of the normalized fission spectrum.

\textbf{Maxwellian average.}
Extend the ideal Maxwellian to all $E>0$:
\[
 w_T(E)=\frac{E e^{-E/\theta}}{\theta^2},
 \qquad \int_0^\infty E e^{-E/\theta}\dd E=\theta^2.
\]
For a $1/v$ cross section,
$\sigma_x(E)=\sigma_0\sqrt{E_{\mathrm{ref}}/E}$,
\begin{boardbox}{Average a $1/v$ absorber}
\begin{align*}
 \bar\sigma_{x,T}
 &=\frac{\sigma_0\sqrt{E_{\mathrm{ref}}}}{\theta^2}
   \int_0^\infty E^{1/2}e^{-E/\theta}\dd E\\
 &=\MathBlank[1.60in]{\frac{\sqrt\pi}{2}\sigma_0
   \sqrt{E_{\mathrm{ref}}/\theta}}.
\end{align*}
Use $\int_0^\infty E^{1/2}e^{-E/\theta}\dd E
=(\sqrt\pi/2)\theta^{3/2}$.
\end{boardbox}
\textbf{Resonance integral.}
For $1/E$ weighting on $[E_\ell,E_u]$,
\[
 I_x(E_\ell,E_u)=\int_{E_\ell}^{E_u}\frac{\sigma_x(E)}{E}\dd E,
 \qquad
 \boxed{\bar\sigma_{x,I}=\frac{I_x}{\ln(E_u/E_\ell)}}.
\]
$I_x$ has cross-section units. Its bounds and any subtraction
of a $1/v$ contribution must match the data definition.
\begin{takeawaybox}[Do not mix weighting assumptions]
The NR average uses the weight $1/[\Sigma_t(E)E]$.
An ordinary $1/E$ resonance integral does \emph{not} include
that flux depression.
\end{takeawaybox}

\HandoutSection{6}{Source-driven water tank}
Treat Prof.~McNeil's water tank as infinite and homogeneous.
Trace ${}^{252}$Cf supplies $s=10\unit{cm^{-3}\,s^{-1}}$;
ignore neutron interactions with the trace source. Assume
$\phi_T=\phi_I=\phi_F$.
For water molecule density $N_w$,
\[
 \bar\Sigma_a=
 \frac{N_w}{3}\sum_{j=T,I,F}
 \left(2\bar\sigma^{\mathrm H}_{a,j}+\bar\sigma^{\mathrm O}_{a,j}\right).
\]
Use microscopic cross sections in $\mathrm{cm^2}$
($1\unit{b}=10^{-24}\unit{cm^2}$). Find
\[
 \phi_{\mathrm{tot}}=\MathBlank[1.00in]{10/\bar\Sigma_a},
 \qquad \phi_j=\MathBlank[0.83in]{\phi_{\mathrm{tot}}/3}.
\]
\SolutionCue{The slide deck's rough estimate
$\bar\Sigma_a\approx0.007\unit{cm^{-1}}$ gives
$\phi_{\mathrm{tot}}\approx1.4\times10^3\unit{cm^{-2}\,s^{-1}}$.
}

\switchcolumn
\RaggedRight
\HandoutSection{7}{Source-free multiplication}
With no external source, physical steady-state balance is
\[
 \bar\Sigma_a\phi_{\mathrm{tot}}
 =\overline{\nu\Sigma_f}\phi_{\mathrm{tot}}.
\]
For a \emph{nonzero} flux, production must equal absorption.
This is the critical condition, but it does not determine the flux magnitude.
\begin{fixedbox}{Introduce the infinite-medium eigenvalue}
Scale only fission production to obtain a stationary eigenproblem:
\[
 \bar\Sigma_a\phi_{\mathrm{tot}}
 =\frac{1}{k_{\infty}}\overline{\nu\Sigma_f}\phi_{\mathrm{tot}},
 \qquad
 \boxed{k_{\infty}=\frac{\overline{\nu\Sigma_f}}{\bar\Sigma_a}}.
\]
\end{fixedbox}
With the consistent eigen-spectrum this is the eigenvalue;
a prescribed approximate spectrum gives an estimate.
\begin{checkpoint}[Interpret the ratio]
Classify the infinite system:
\begin{align*}
 k_{\infty}<1 &: \quad\RevealBlank[1.25in]{subcritical},\\
 k_{\infty}=1 &: \quad\RevealBlank[1.25in]{critical},\\
 k_{\infty}>1 &: \quad\RevealBlank[1.25in]{supercritical}.
\end{align*}
Does doubling $\phi(E)$ change $k_{\infty}$? Does it change
$R_f$?
\RuledLines{2}
\end{checkpoint}

\HandoutSection{8}{U-235 dissolved in water}
Replace the external source by dissolved, pure ${}^{235}$U.
Ignore other solute constituents and leakage. Define
\[
 r=\frac{N_{\mathrm U}}{N_w},
 \qquad \phi_T=\phi_I=\phi_F.
\]
For each nuclide, average over the \emph{same mixture spectrum};
with equal range fluxes, $\bar\sigma_x^i=\frac13\sum_j\bar\sigma_{x,j}^i$.
Only uranium contributes fission production in this model:
\begin{align*}
 \bar\Sigma_a
 &=N_w\left(2\bar\sigma_a^{\mathrm H}
     +\bar\sigma_a^{\mathrm O}+r\bar\sigma_a^{\mathrm U}\right),\\
 \overline{\nu\Sigma_f}
 &=N_wr\,\overline{\nu\sigma_f}^{\mathrm U}.
\end{align*}
\begin{boardbox}{Form production divided by absorption}
\[
 \boxed{k_{\infty}\approx
 \frac{r\,\overline{\nu\sigma_f}^{\mathrm U}}
 {2\bar\sigma_a^{\mathrm H}+\bar\sigma_a^{\mathrm O}
      +r\bar\sigma_a^{\mathrm U}}}.
\]
Why does $N_w$ cancel? What is being held fixed when this
expression is used to vary $r$?
\RuledLines{2}
\end{boardbox}
\begin{takeawaybox}
Changing composition generally changes the spectrum and its
averages. Equal thermal, intermediate, and fast fluxes are an
\emph{assumption for this example}, not a prediction.
\end{takeawaybox}
\end{handoutcolumns}
\end{document}
